\begin{exercise}[subtitle={Improvement of the subgaussian bound \Coffeecup\Coffeecup}]
    \label{ex: subgaussian bound improvement}
    In the table of Theorem \ref{thm: subgaussian equivalences} prove
    \(c_{42} \le \sqrt{3}\). That is
    \[
        K_4 \le \sqrt{3}K_2.
    \]
    \textbf{Hint:} Skip the bounds on the Gamma function  that lead from \(K_2\)
    to \(K_2'\) in the proof of Theorem \ref{thm: subgaussian equivalences} 
    and directly move to ``\ref{it: subgaussian moment decay} \(\Rightarrow\)
    \ref{it: subgaussian X^2 sub-exponential}'', where the Stirling
    approximation step becomes unnecessary to get
    \[
        \E{\exp(t^2 X^2)} \le \frac{1+(t K_2)^2}{1-(t K_2)^2}.
    \]
    Then consider the definition of \(K_4\).
\end{exercise}

\begin{solution}
    By definition we have \(\E{\abs{X}^p} \le K_2^p p \Gamma(\frac{p}2)\) and 
    therefore  
    \[
        \E{\abs{X}^{2p}} \le K_2^{2p} (2p) \Gamma(p) = 2K_2^{2p} p!
    \]
    consequently
    \begin{align*}
        \E{\exp(t^2 X^2)}
        &= \sum_{p=0}^\infty \frac{t^{2p} \E{\abs{X}^{2p}}}{p!}
        \\
        &\le 1+ \sum_{p=1}^\infty \frac{2(t K_2)^{2p} p!}{p!}
        \overset{\substack{\text{geom.}\\\text{series}}}= 1 + 2\frac{(t K_2)^2}{1-(t K_2)^2}
        = \frac{1+(t K_2)^2}{1-(t K_2)^2}.
    \end{align*}
    Observe that
    \[
        \frac{1+\lambda^2}{1-\lambda^2} \overset{!}\le 2
    \]
    is satisfied for \(\lambda^2 \le \frac13\) or equivalently \(t^2 \le \frac{1}{3K_2^2}\) for the choice of \(\lambda = tK_2\).
    Choosing \(t=\frac1{c}\) we get consequently get for \(\frac1{c^2} \le \frac{1}{3K_2^2}\)
    \[
        \E{\exp(\tfrac{X^2}{c^2})} \le 2.
    \]
    Since this is guaranteed to hold for all \(c \ge \sqrt{3}K_2\) and \(K_4\)
    may be chosen as the smallest constant where this holds, we have \(K_4 \le
    \sqrt{3}K_2\).
\end{solution}
